\section{Du modèle aux flux numériques}
\label{sec:schemes}

\subsection{Grille, indices et trois niveaux de description}

Une image numérique possède $N_y$ lignes et $N_x$ colonnes. L'indice $i$
désigne la direction verticale, de pas $\dy$, et $j$ la direction
horizontale, de pas $\dx$. On peut interpréter les valeurs $U_{ij}$
comme celles de centres de cellules
\begin{equation}
 x_j=(j+\tfrac12)\dx,\qquad y_i=(i+\tfrac12)\dy,
 \quad 0\leq j<N_x,\quad 0\leq i<N_y,
 \label{eq:cell-centres}
\end{equation}
sur $\Omega=(0,N_x\dx)\times(0,N_y\dy)$. Dans les expériences,
$N_x=N_y=128$ et $\dx=\dy=1$ pixel. La matrice est convertie en
\texttt{float64} et copiée avant toute évolution. Les données doivent
être finies, à deux dimensions et comporter au moins deux pixels dans
chaque direction.

Il convient de séparer trois objets. L'EDP continue décrit une fonction
de l'espace et du temps. Le système \emph{semi-discret} est une équation
différentielle sur les $N_xN_y$ valeurs, après remplacement des
dérivées spatiales par des échanges sur la grille. Le schéma
\emph{entièrement discret} remplace ensuite cette équation différentielle
par une suite d'états $U^n$ au temps $t_n=n\dt$. Les propriétés de l'un
ne se transmettent pas automatiquement aux deux autres.

\subsection{Une conductance unique par interface intérieure}

Entre deux pixels horizontalement voisins, on définit
\begin{equation}
 \delta^x_{i,j+1/2}=U_{i,j+1}-U_{ij},\qquad
 a^x_{i,j+1/2}=c\left(\frac{|\delta^x_{i,j+1/2}|}{\dx}\right).
 \label{eq:horizontal-interface}
\end{equation}
Les interfaces verticales sont décrites de même par
\begin{equation}
 \delta^y_{i+1/2,j}=U_{i+1,j}-U_{ij},\qquad
 a^y_{i+1/2,j}=c\left(\frac{|\delta^y_{i+1/2,j}|}{\dy}\right).
 \label{eq:vertical-interface}
\end{equation}
Pour la chaleur, toutes les conductances intérieures valent un.
Pour Perona--Malik, on utilise l'une des deux lois
\eqref{eq:pm-conductances}. La conductance d'une interface est
partagée par ses deux pixels ; on ne calcule pas deux coefficients
indépendants selon le sens d'échange.

Le flux de gradient orienté dans le sens des indices croissants est
$F^x=a^x\delta^x/\dx$, et de même $F^y=a^y\delta^y/\dy$.
Le flux physique de diffusion est son opposé. Le code stocke
directement les contributions à la divergence,
\begin{equation}
 Q^x_{i,j+1/2}=\frac{a^x_{i,j+1/2}\delta^x_{i,j+1/2}}{\dx^2},
 \qquad
 Q^y_{i+1/2,j}=\frac{a^y_{i+1/2,j}\delta^y_{i+1/2,j}}{\dy^2}.
 \label{eq:edge-contributions}
\end{equation}
Une contribution $Q^x$ est ajoutée au pixel gauche et soustraite
au pixel droit ; $Q^y$ est ajoutée au pixel supérieur et soustraite
au pixel inférieur. Par exemple, si le pixel droit est plus clair,
$Q^x>0$ : le pixel gauche augmente et le droit diminue. Cette convention
de signe réalise le lissage tout en conservant la somme.

\subsection{Frontières et actualisation d'Euler}

Seules les interfaces entre pixels du domaine sont créées. On pose
les contributions des interfaces extérieures égales à zéro :
\begin{equation}
 Q^x_{i,-1/2}=Q^x_{i,N_x-1/2}=0,\qquad
 Q^y_{-1/2,j}=Q^y_{N_y-1/2,j}=0.
 \label{eq:boundary-edges}
\end{equation}
Il n'existe donc ni raccordement périodique ni échange avec une valeur
extérieure. Pour la chaleur, cela revient à utiliser un pixel fantôme
égal au pixel frontière, dans l'interprétation cellulaire
\eqref{eq:cell-centres}. Un pixel de coin possède deux voisins,
un pixel de bord non situé dans un coin en possède trois, et un pixel
intérieur en possède quatre.

Le système semi-discret est
\begin{equation}
 \frac{\dd U_{ij}}{\dd t}=(L_h(U))_{ij}
 =Q^x_{i,j+1/2}-Q^x_{i,j-1/2}
  +Q^y_{i+1/2,j}-Q^y_{i-1/2,j}.
 \label{eq:semidiscrete-directional}
\end{equation}
Euler explicite donne
\begin{equation}
 U^{n+1}_{ij}=U^n_{ij}+\dt(L_h(U^n))_{ij}.
 \label{eq:explicit-euler}
\end{equation}
Toutes les différences et toutes les conductances proviennent de
$U^n$. Le tableau $U^{n+1}$ est distinct ; une actualisation pixel
par pixel en place mélangerait deux niveaux temporels et ne serait
plus ce schéma. Pour un pixel intérieur, la chaleur est explicitement
\begin{align}
 U^{n+1}_{ij}=U^n_{ij}
 &+\frac{\dt}{\dx^2}
    (U^n_{i,j+1}-2U^n_{ij}+U^n_{i,j-1})\notag\\
 &+\frac{\dt}{\dy^2}
    (U^n_{i+1,j}-2U^n_{ij}+U^n_{i-1,j}).
 \label{eq:heat-stencil}
\end{align}
Sur le bord, les termes correspondant à un voisin extérieur sont
simplement absents dans la forme conservative.

\subsection{Un algorithme et des sorties fidèles à l'évolution}

Le pseudocode suivant expose les opérations essentielles. Les signes
opposés sont appliqués à une divergence initialement nulle, non au
tableau courant pendant le calcul des flux.

\begin{lstlisting}[basicstyle=\ttfamily\small,frame=single]
U = finite_float64_copy(f)
save_copy(U, n=0) if 0 is requested
for n in 1, ..., largest_checkpoint:
    delta_x = U[:, 1:] - U[:, :-1]
    delta_y = U[1:, :] - U[:-1, :]
    a_x = c(abs(delta_x)/dx)   # or ones for heat
    a_y = c(abs(delta_y)/dy)
    q_x = a_x * delta_x / dx**2
    q_y = a_y * delta_y / dy**2
    D = zeros_like(U)
    D[:, :-1] += q_x; D[:, 1:] -= q_x
    D[:-1, :] += q_y; D[1:, :] -= q_y
    U_new = U + dt*D
    U = U_new
    save_copy(U, n) if n is requested
\end{lstlisting}

Une trajectoire est poursuivie jusqu'au plus grand checkpoint demandé.
Les états intermédiaires sont copiés, et non recalculés en repartant
de l'observation. Cela assure une comparaison à budget d'itérations
commun. Le checkpoint zéro est une copie de la donnée initiale pour
chaque trajectoire ; le CSV conserve ces lignes par méthode et une
ligne bruitée commune, sans générer de nouvelles observations.
Aucune étape ne borne les intensités dans $[0,1]$.

Le chronométrage des solveurs commence après la copie initiale et
l'éventuelle sauvegarde à $n=0$. Il comprend les étapes d'évolution
et les copies des checkpoints positifs. Les durées enregistrées sont
cumulatives depuis ce départ ; elles ne doivent donc pas être
additionnées entre checkpoints. Elles n'incluent ni création du
bruit, ni calcul des métriques, ni figures, ni sérialisation finale.

\subsection{Portée de la variante directionnelle}

Une différence $|\delta^x|/\dx$ approche $|\partial_xu|$ sur une
interface horizontale, non $\sqrt{(\partial_xu)^2+(\partial_yu)^2}$.
Même pour une image régulière, les deux coefficients évalués aux
interfaces horizontales et verticales peuvent donc être différents.
À grille raffinée et pour une fonction suffisamment régulière,
la forme directionnelle correspond formellement, à l'intérieur,
à l'opérateur
\begin{equation}
 \partial_x\bigl(c(|\partial_xu|)\partial_xu\bigr)
 +\partial_y\bigl(c(|\partial_yu|)\partial_yu\bigr),
 \label{eq:directional-limit}
\end{equation}
et non à l'opérateur scalaire isotrope de
\eqref{eq:pm-continuous}. Cette différence ne disparaît pas en
diminuant seulement $\dt$. Une rotation de l'image par rapport aux
axes peut modifier les conductances et les résultats. La variante
est simple, conservative et courante pour une première étude,
mais cette dépendance directionnelle est une limite intrinsèque
de l'objet numérique évalué.
